% Laget automatisk av elevark/lag_elevark.py fra opplegg/g-med-mobil.html. Endringer her blir overskrevet.
\documentclass[11pt]{article}
\usepackage{elevark}
\begin{document}
\elevarkhode{Fysikk 1 · Rettlinjet bevegelse}{Mål $g$ med mobilen}{https://fysikklab.no/opplegg/g-med-mobil.html}{Opplegget på nett}

En stålkule dyttes ut over kanten av et bord og faller i gulvet. Mobilen lytter: Klikket når kula dyttes, starter klokka, og smellet i gulvet stopper den. Med falltiden fra flere høyder kan dere bestemme tyngdeakselerasjonen $g$.

\section*{Hensikt}

\begin{itemize}
\item Bestemme $g$ fra falltiden for flere høyder.
\item Linearisere $h = \tfrac{1}{2}gt^{2}$ ved å plotte $h$ mot $t^{2}$, og bruke stigningstallet.
\end{itemize}

\section*{Utstyr}

\begin{itemize}
\item Mobil med appen phyphox (gratis) og eksperimentet «Akustisk stoppeklokke» («Acoustic Stopwatch»). Det starter på én lyd og stopper på den neste.
\item En liten stålkule og en metallinjal. Et hardt gulv, et bord og gjerne en stol, en hylle eller en trappeavsats for andre høyder. Målebånd.
\end{itemize}

\section*{Teorien}

Kula faller fra ro, så $h = \tfrac{1}{2}gt^{2}$. Horisontal startfart endrer ikke falltiden. Grafen av $h$ mot $t^{2}$ er en rett linje med stigningstall $g/2$, og hver måling gir også $g = 2h/t^{2}$.

\section*{Framgangsmåte}

\begin{enumerate}
\item Legg kula helt ute på kanten av bordet. Mål høyden $h$ fra underkanten av kula til gulvet.
\item Legg mobilen på bordet like ved kula. Start den akustiske stoppeklokka, og juster terskelen slik at vanlig prat ikke starter den.
\item Slå lett på kula sidelengs med linjalen, slik at den ruller ut over kanten med lav fart. Klikket starter klokka, og smellet i gulvet stopper den.
\item Gjør tre målinger fra hver høyde og bruk gjennomsnittet. Gjenta for minst fem høyder mellom 0,4~m og 1,5~m.
\end{enumerate}

\section*{Oppgaver}

\begin{oppgaveliste}
\oppgave \fase{gjett}{Gjett}Hvor lang tid tror dere det tar for kula å falle 1~m? Hvor mye lengre tid tar det fra 2~m?
\svar
\oppgave \fase{test}{Mål}Gjør målingene, eller trykk på «Fyll inn eksempel». Hvilken verdi får dere for $g$?
\svar
\oppgave \fase{forklar}{Forklar}Hvorfor plotter vi $h$ mot $t^{2}$ og ikke mot $t$? Hvorfor er stigningstallet $g/2$?
\svar
\oppgave \fase{test}{Regn}Lyden går med omtrent 343~m/s. Mobilen ligger på bordet, så smellet fra gulvet bruker litt tid før det når mobilen. Hvor stor blir den forsinkelsen fra 1,0~m, og hvor mye endrer den $g$?
\svar
\oppgave \fase{forklar}{Vurder}Sammenlign $g$ fra stigningstallet med gjennomsnittet av enkeltmålingene. Hvilken metode stoler dere mest på? Hvordan kan dere plassere mobilen for å gjøre feilen fra lydforsinkelsen mindre?
\svar
\oppgave \fase{utfordring}{Utfordring}Gjenta med en pingpongball fra den største høyden. Får dere samme $g$? Hva forteller det?
\svar
\end{oppgaveliste}

\end{document}
